
Mamba-2 establishes theoretical equivalence between Transformer attention and state space models, suggesting a principled path for knowledge transfer from pretrained Transformers to efficient SSM architectures. This work tests whether this duality enables practical distillation by extracting SSM parameters directly from BERT attention weights. The experiments reveal that while duality-derived parameters are numerically stable (0\% NaN/Inf, 0.$11\times$ magnitude ratio relative to Transformer outputs), they reconstruct attention outputs 2.13\% worse than random initialization. This negative effect is consistent across all 12 BERT layers (p < 0.0001, Cohen's d = $-4.56)$. Two probable causes are identified: the output-level Frobenius norm metric may not capture what duality preserves, and BERT's multi-head attention may violate the theoretical SSD conditions requiring scalar A matrices and 1-semiseparable causal masks. These results provide empirical scope clarification for Mamba-2 duality: valid parameters do not imply preserved structure. Practitioners considering attention-to-SSM conversion should test reconstruction quality, not just parameter validity.
